%%%%%%%%%%%%%%%%%%%%%%%%%%%%%%%%%%%%%%%%%%%%%%%%%%%%%%%%%%%%%%%%%%%%%%%%%%%
\chapter{Introduction}
%%%%%%%%%%%%%%%%%%%%%%%%%%%%%%%%%%%%%%%%%%%%%%%%%%%%%%%%%%%%%%%%%%%%%%%%%%%


Over the last century, humankind has finally set its first step in the expanse of space and acquired vast amounts of data, 
providing massive boons in multiple fields and new technologies that we take for granted today, which were almost fantastical 
at the turn of the 20th century, such as the GPS. There are great incentives in continuing to invest in space ventures, 
as the future may bring lunar colonies, space elevators, asteroid mining, and many new benefits. One way to achieve this 
is through educating the younger generation about astronomy, which is currently being done by educational institutions 
globally to varying degrees of depth and knowledge. However,  traditional astronomy education often glosses over basic 
information about space, planets, and other celestial objects, which leads to a lack of sustained engagement in the field 
of astronomy. Students often stop learning about space by senior high school in a K-12 program. Additionally, despite the 
abundance of resources and learning tools about astronomy, many of these programmes are intended for learners with a good 
foundation of astronomical knowledge and a self-sustained interest in the subject matter. The planned application aims to 
enhance the learning experience for astronomy students by implementing an interactive game-based learning system for space 
exploration, with emphasis on collaboration. A learning system focused on a fun feedback loop may be effective in capturing 
 interests of students by being challenged to learn astronomical concepts and collaborating with other learners to collectively
  progress their knowledge.



\section{Project Context}

Aphelion is a web-based game designed to provide an effective entry point for students who are interested in learning astronomy. 
The application prioritizes its game and collaboration aspects, ensuring that users learn and train their concepts while having fun. 
Students should be immersed and feel like they are exploring the stars, with the actual learning happening in the background.

The game features a 2D starmap divided into sectors. Each star sector contains 50-70 star systems handpicked from the HYG Database, 
based on their brightness and proximity to other stars. In the beginning, every star system is “unexplored”, with one star system 
preemptively explored to serve as the home base of users. From this home base, neighbouring stars are available to begin an exploration 
mission. This mission forms the core of the collaboration aspect – to fully map out a star system, users must progress the exploration 
meter by accomplishing research minigames. These minigames test the user’s knowledge of basic astronomical concepts and range from 
guessing the Goldilocks zone of a star to sending probes to its celestial bodies, among other tasks, which are randomized every time 
the user attempts a research minigame. Once a star system is fully explored, it becomes a point of exploration, allowing users to read 
about the star system and its planets with detailed text descriptions of their features, landscapes, and other facts if available, and 
further systems nearby can have their expedition initiated. The provided information of the stars and planets will be interpreted from 
NASA’s database.

The reason for this project’s initiation comes from a lack of easily accessible and game-based learning tools to serve as a gateway for 
students who are interested in learning astronomy or pursuing astronomy and space exploration careers. There are studies that showcase 
an increased effectiveness of learning using educational games [28]. Astronomy was chosen to be provided with a game-based learning tool 
due to its upcoming importance in future research and industry [7].


\section{Purpose and Description}

The project aims to integrate complex astronomical data into beginner-level science education by creating a game-based educational 
platform. Given the vast number of educational resources, only a few cater to beginner learners. Specifically, an educational platform 
for astronomy with a well-defined and engaging core loop. This project will utilize gamification mechanics, such as progress tracking 
and collaborative elements, to promote motivation and knowledge retention among students by employing publicly accessible and credible 
data.

This project is a web application that incorporates interactive models, creating an immersive learning experience for students. 
The platform will collect data from the NASA Exoplanet API and HYG Database as sources. It will be built upon the instructional 
structure of the ADDIE (Analyze, Design, Develop, Implement, Evaluate) framework, as it is a highly effective and reliable model 
[22][23]. The system distinguishes itself through content gamification, which integrates game design elements using astronomical data. 
The platform will feature engaging exploration through a 2D starmap, collaborative learning, and progressive feedback loops.


\section{Objectives}

The main objective of this project is to create a game-based educational website on astronomy. Students will collaborate with 
each other to simulate the exploration of a star system and its planets. By creating a game-based learning environment, the interest 
of students in the subject of astronomy, as well as their efficacy in learning, may increase.

\section{Scope and Limitations}

The Aphelion project is primarily an educational tool for K-12 students. Its focus is to simulate collaboration and space exploration in 
a 2D star map. The project will only operate on the content integration of data from the selected sources and the application of the 
ADDIE model as instructional design for learning. Furthermore, the project will be limited to 2D graphics to ensure accessibility across 
various web browsers and hardware. Complex astrophysical formulas are simplified into interactive mechanics rather than raw mathematical 
data. Lastly, the application is only available with a stable internet connection for data retrieval and collaborative synchronization. 
